\section{Capítulo 1}
\subsection{Introducción}
El kéfir (del turco $kief$, que significa ``sentirse bien'' o ``agradable sensación'') es una
bebida fermentada a partir de leche y un cultivo madre de microorganismos
presentes en los granos de kéfir, conocidos en México como búlgaros. Esta bebida
representa una de las prácticas de fermentación artesanal más antiguas,
considerada incluso uno de los primeros productos pecuarios utilizados por el ser
humano, con orígenes en las montañas del Cáucaso y Europa del Este \cite{KumarBehera2017} \cite{LucianaGabriela;2019}.
Con el paso del tiempo, su producción y consumo se han extendido hacia diversas
regiones del mundo, incluyendo México, donde fue adoptada tanto por sus
propiedades nutricionales como por su valor probiótico. Además, ha sido adaptada
por productores artesanales que han incorporado métodos propios, influenciados
por la diversidad cultural y gastronómica del país \cite{CamachoVera2019}  \cite{CesarOjeda2021}.\newline
En la actualidad, la elaboración artesanal de kéfir en México ha cobrado gran
relevancia debido a su creciente demanda comercial y al interés científico, lo cual
ha derivado en una variedad de métodos de producción \cite{MauroCadme2021}. Estos métodos se ven
condicionados por múltiples factores, como el tipo de leche utilizada, la proporción
de búlgaros, la temperatura, el tiempo de fermentación y la fuente del inóculo (es
decir, el equipo o recipiente utilizado) \cite{LopezRuiz2023}.\newline
Siendo uno de los aspectos más importantes en el proceso de fermentación es la
recolección del kéfir en el momento adecuado, ya que el tiempo de fermentación
influye directamente en el sabor de la bebida. Sin embargo, en base a los
ingredientes usados varía el punto óptimo de recolección lo que añade
incertidumbre al proceso, obligando a los pequeños y medianos productores a
definir sus propios parámetros iniciales para sus fermentaciones mediante prueba y
error para definir el tiempo adecuado de recolección del kéfir. Esto no solo genera
una gran variabilidad en la calidad del producto final, sino que también implica
pérdidas económicas significativas, ya que lotes completos pueden resultar
inadecuados para el consumo o la venta \cite{RuedaManzano1980}.\newline
Por esta razón el monitoreo y control de los parámetros del proceso de fermentación
se vuelve fundamental para garantizar una producción adecuada del kéfir, de lo
contrario se obtendría una bebida avinagrada o muy ácida, siendo no apta para el
consumo humano. Con el fin de atender esta problemática, el presente proyecto
propone el diseño de un prototipo de monitoreo para la fermentación artesanal del
kéfir, proporcionando una herramienta tecnológica que ayude a determinar el tiempo
adecuado para la recolección del kéfir. De esta manera, contribuir a la mejora de la
producción artesanal de kéfir, ofreciendo a los productores un mayor monitoreo
sobre su proceso de producción, garantizando la calidad del producto final.